\section{Introduction}
\label{sec:introduction}

\subsection{Musculoskeletal deconditioning and long-duration missions}

Skeletal muscle adapts to the load it carries. When the load is removed---in weightlessness, or when a healthy
person lies in bed for weeks---the muscles that normally hold the body up against gravity lose mass and strength.
For a crew on a transit to Mars, which lasts on the order of six months, that loss is not an abstraction: a crew
has to be able to stand, move and work on arrival, and countermeasures have to be planned before departure. A
planner therefore needs two numbers that the literature does not currently provide in one place: how much of a
given muscle is lost after a given number of days of unloading, and how much that depends on anything other than
time and the muscle in question.

\subsection{Bed rest as a terrestrial analogue}

Actual spaceflight data on lower-limb muscle size are scarce, and the few available cohorts are small. The
established terrestrial analogue is bed rest: healthy volunteers remain in bed, usually tilted six degrees
head-down so that fluid shifts towards the head as it does in orbit, for days to months, while their muscles are
scanned before, during and after the unloading. Related models unload one leg (unilateral lower-limb suspension)
or the whole body in water (dry immersion). Over three decades these campaigns---run at NASA, MEDES in Toulouse,
the Charit\'e in Berlin, the DLR in Cologne, and elsewhere---have produced a large and heterogeneous literature
on muscle loss~\citep{leblanc1992,alkner2004,belavy2009,miokovic2012,trappe2007,trappe2024sprint,dulac2024}.

\subsection{The gap}

Each campaign reports its own muscles, at its own time points, with its own imaging protocol, participants and
countermeasures. A reader who wants to know how much soleus is left after sixty days has to read dozens of papers
and reconcile them by hand. Two further difficulties make naive pooling misleading. First, many papers report the
same participants: a single campaign can yield five or more publications, so counting papers overstates the
evidence and, worse, leaks information across the folds of any cross-validation that treats papers as
independent. Second, the numbers themselves are heterogeneous: some papers report muscle volume by magnetic
resonance imaging (MRI), others cross-sectional area by computed tomography (CT), lean mass by dual-energy X-ray
absorptiometry (DXA), or thickness by ultrasound; some report a composite such as the quadriceps together with its
four heads; and some report only a percentage change.

\subsection{Aim and contributions}

The project had one aim, taken from the accepted abstract: to develop a literature-derived dataset and a
machine-learning framework for predicting lower-limb muscle atrophy in spaceflight analogues. It makes four
contributions.

\begin{enumerate}
  \item \textbf{A frozen, documented dataset.} 742 comparisons against baseline from 52 studies and 36 independent
        campaigns, each traceable to a table, figure or page of its source, extracted into a schema frozen before
        extraction began and released under a hash-verified version tag (\cref{sec:methods-data,sec:results-data}).
  \item \textbf{A campaign map.} Every study is assigned to the campaign that produced its participants, which
        makes the campaign---not the paper---the unit of validation (\cref{sec:cohort-map}).
  \item \textbf{A three-tier modelling framework} in which a three-level meta-regression carries the scientific
        estimate, four machine-learning families (and, post hoc, a pretrained tabular neural network) test whether
        flexible models find structure the curve misses, and
        a probabilistic language-model forecast tests whether a model that reads each campaign's description can do
        what the numeric columns cannot (\cref{sec:methods-models}).
  \item \textbf{A declared evaluation.} Every threshold, reading rule and sensitivity analysis was written down
        before the answer it governs, and every number in this report is regenerated from the frozen dataset by
        one command (\cref{sec:reproducibility}).
\end{enumerate}

The report is written as a framework-and-dataset contribution, not a performance contribution. With 32 independent
campaigns in the modelling subset, no cross-validated error will impress a machine-learning audience, and it does
not need to: the questions that matter here are whether the estimates are physiologically plausible, whether the
uncertainty is stated honestly, and whether the approach could ever inform mission planning.

\subsection{Structure of the report}

\Cref{sec:background} summarises the physiological and methodological background. \Cref{sec:methods-data}
describes how the dataset was built and \cref{sec:methods-models} how it was modelled. \Cref{sec:results-data}
describes the corpus and the dataset, and \cref{sec:results-models} reports the three tiers.
\Cref{sec:discussion} interprets the findings, states the limitations in full, lists the work that remains and
reconciles the results with the accepted abstract. The appendices hold the verbatim search strategies, the
extraction schema, the campaign table, the muscle classification and the full model outputs.
